\chapter{安装 NixOS 虚拟机}
\chapterlead{本章在一台现有计算机上安装 NixOS 虚拟机，完成后可登录、联网、修改配置并回滚。虚拟机的代价是可随时删除重建。后续章节的实验均在这台机器上完成。}

\section{先决条件}

宿主机需满足：内存 8\,GB 以上，可用磁盘 30\,GB 以上。虚拟机分配 2 个处理器核心、4\,GB 内存、20\,GB 磁盘，足够本书使用。

虚拟化软件按宿主系统选择。

\begin{booktable}{\textwidth}{L{0.980} L{1.194} J{0.826}}
\toprule
软件 & 平台 & 说明 \\
\midrule
VirtualBox & Windows、Linux、Intel Mac & 免费，界面直观，本章以此为例 \\
GNOME Boxes & Linux & 多数发行版可直接安装，配置最少 \\
UTM & Apple Silicon Mac & 免费，支持 aarch64 虚拟机 \\
VMware Workstation & Windows、Linux & 性能较好，个人使用免费 \\
\bottomrule
\end{booktable}

Apple Silicon 的 Mac 需下载 \texttt{aarch64} 镜像，其余平台下载 \texttt{x86\_64} 镜像，安装步骤相同。

\section{下载镜像}

访问 \url{https://nixos.org/download}，选择 NixOS 26.11 的图形安装镜像（GNOME ISO），约 2 至 3\,GB。镜像分带图形界面与最小化两个版本；本书使用前者，安装过程中可在图形环境里查资料。

下载后核对校验值。SHA-256 以官网公布的值为准，不从第三方博客抄录。

\begin{shellsession}
$ sha256sum nixos-gnome-26.11-x86_64-linux.iso
\end{shellsession}

Windows 使用 PowerShell 的 \texttt{Get-FileHash}，macOS 使用 \texttt{shasum -a 256}。输出与官网一致后方可继续。

\section{创建虚拟机}

以 VirtualBox 为例。新建虚拟机，类型 Linux，版本 Other Linux (64-bit)；内存 4096\,MB，处理器 2 个，磁盘 20\,GB。在“系统”设置中启用 EFI，在“显示”中将显存调至 128\,MB。

创建完成后，在“存储”中把光盘指向下载的 ISO，启动。

VirtualBox 有一处常见问题：默认图形控制器可能导致安装界面显示异常。启动后若黑屏，关机，将“显示 \textrightarrow\ 图形控制器”改为 VMSVGA，再次启动。

\begin{notebox}[为什么启用 EFI]
NixOS 默认使用 systemd-boot 引导，需要 EFI 分区。虚拟机启用 EFI，实体机以 UEFI 模式启动安装介质，两者保持一致，配置可以直接复用。
\end{notebox}

\section{进入安装环境}

启动后进入 GNOME 桌面，桌面上有“Install NixOS”图标与一个终端。图形安装器可以完成基本安装，本书采用命令行方式，原因有两点：命令可复制、可比对、可记入笔记；图形安装器的每个选项最终也写入同一份配置文件。

打开终端，确认网络与磁盘。

\begin{shellsession}
$ ping -c 3 nixos.org
$ lsblk
NAME   MAJ:MIN RM  SIZE RO TYPE MOUNTPOINTS
sr0     11:0    1  2.4G  0 rom
vda    254:0    1   20G  0 disk
\end{shellsession}

\texttt{vda} 是虚拟磁盘。名称中的 \texttt{a} 表示第一块盘，分区后在名称后追加序号，成为 \texttt{vda1}、\texttt{vda2}。实体机通常是 \texttt{nvme0n1} 或 \texttt{sda}，第十八章说明如何确认。

\section{分区与格式化}

下列命令会清空整块 \texttt{vda}。虚拟机中无数据损失，但应养成固定习惯：执行前先用 \texttt{lsblk} 确认磁盘名。

\begin{shellsession}
$ sudo parted /dev/vda -- mklabel gpt
$ sudo parted /dev/vda -- mkpart ESP fat32 1MiB 512MiB
$ sudo parted /dev/vda -- set 1 esp on
$ sudo parted /dev/vda -- mkpart root ext4 512MiB 100%
$ sudo mkfs.fat -F 32 -n boot /dev/vda1
$ sudo mkfs.ext4 -L nixos /dev/vda2
$ sudo mount /dev/disk/by-label/nixos /mnt
$ sudo mkdir -p /mnt/boot
$ sudo mount /dev/disk/by-label/boot /mnt/boot
\end{shellsession}

分区只有两个：512\,MB 的 EFI 分区存放引导程序，其余空间全部给根分区。该布局简单，覆盖本书全部实验。需要交换空间、独立 \texttt{/home} 或磁盘加密的情形，在第十八章的基础上调整。

\texttt{mkfs} 意为“建立文件系统”。刚完成分区的磁盘仅是一段空闲区域，\texttt{mkfs.ext4} 在其上写入目录结构、空闲块表等元数据，操作系统才能存放文件。EFI 分区必须使用 FAT32，这是固件的要求。

\section{生成初始配置}

\begin{shellsession}
$ sudo nixos-generate-config --root /mnt
writing /mnt/etc/nixos/hardware-configuration.nix...
writing /mnt/etc/nixos/configuration.nix...
\end{shellsession}

命令生成两份文件。\texttt{hardware-\allowbreak{}configuration.\allowbreak{}nix} 记录磁盘、分区、内核模块与 CPU 特性，不应手改；重装或更换硬件时重新生成。\texttt{configuration.\allowbreak{}nix} 是需要编辑的文件，初始状态多为注释。

用 nano 打开（首次使用 nano 比 vim 直接：\texttt{Ctrl+O} 保存，\texttt{Ctrl+X} 退出）。

\begin{shellsession}
$ sudo nano /mnt/etc/nixos/configuration.nix
\end{shellsession}

修改下列各项，其余保持默认。

\begin{nixcode}
{ config, pkgs, ... }:

{
  imports = [ ./hardware-configuration.nix ];

  boot.loader.systemd-boot.enable = true;
  boot.loader.efi.canTouchEfiVariables = true;

  networking.hostName = "nixos-lab";
  networking.networkmanager.enable = true;

  time.timeZone = "Asia/Shanghai";
  i18n.defaultLocale = "zh_CN.UTF-8";
\end{nixcode}

文件尚未结束，继续找到用户与软件包部分。

\begin{nixcode}
  services.openssh.enable = true;

  users.users.weiwen = {
    isNormalUser = true;
    extraGroups = [ "wheel" "networkmanager" ];
    initialPassword = "change-me";
  };

  environment.systemPackages = with pkgs; [
    git
    vim
    curl
    htop
  ];

  system.stateVersion = "26.11";
}
\end{nixcode}

各项含义如下。

\begin{itemize}
  \item \texttt{systemd-boot} 是引导程序，负责将内核交给机器。
  \item \texttt{hostName} 是主机名，显示在提示符中。
  \item \texttt{NetworkManager} 管理网络，虚拟机与笔记本均适用。
  \item \texttt{timeZone} 与 \texttt{defaultLocale} 决定时区与语言。\texttt{zh\_CN.UTF-8} 使系统支持中文显示；终端输出建议保持英文，便于检索报错。
  \item \texttt{services.\allowbreak{}openssh} 启用 SSH 服务，第九章使用。
  \item \texttt{extraGroups} 中的 \texttt{wheel} 表示该用户可通过 \texttt{sudo} 取得管理员权限。
  \item \texttt{initialPassword} 是首次登录的临时密码，登录后应立即修改。
  \item \texttt{stateVersion} 标记配置面向的版本，安装后不要修改，原因见第十二章。
\end{itemize}

\begin{warnbox}[缩进与分号]
Nix 以空白分隔，以分号结束每一项，以花括号分组。缺少一个分号会导致整次重建失败。报错会指出行号，按行号定位，不要凭印象猜测。
\end{warnbox}

\section{执行安装}

\begin{shellsession}
$ sudo nixos-install
\end{shellsession}

该命令下载并构建系统，视网络状况需要十几分钟。过程中会要求设置 root 密码。结束时输出重启提示。

\begin{shellsession}
$ sudo reboot
\end{shellsession}

重启后出现引导菜单，默认项即刚安装的系统。登录界面使用创建的用户名与 \texttt{initialPassword} 中的密码。

\section{首次登录后的三项操作}

\begin{shellsession}
$ nixos-version
26.11.20261113.d5cc683 (Zokor)
$ sudo passwd weiwen
$ sudo nixos-rebuild switch
\end{shellsession}

第一项确认版本。第二项替换临时密码。第三项执行一次重建，确认流程可用；此后每次修改配置都使用同一条命令。

\section{练习：改坏并恢复}

本章的核心练习。目标是建立对回滚机制的直接经验：系统配置可以改坏，也有确定的恢复路径。

\begin{enumerate}
  \item 将 \texttt{configuration.\allowbreak{}nix} 中的 \texttt{networking.\allowbreak{}hostName} 改为 \texttt{broken-host}，执行 \texttt{sudo nixos-\allowbreak{}rebuild switch}。退出重新登录，提示符中的主机名随之改变。
  \item 制造一处语法错误，例如删除一个分号，重新构建。记录报错内容，按行号定位问题。
  \item 执行 \texttt{sudo nixos-rebuild switch -{}-rollback} 回到上一代。重新登录，主机名恢复。
  \item 执行 \texttt{sudo nixos-\allowbreak{}rebuild list-\allowbreak{}generations}，统计当前存在几代系统。
\end{enumerate}

\begin{checkpoint}
本章结束时，应能：
\begin{itemize}
  \item 独立创建 NixOS 虚拟机并完成安装；
  \item 说明 EFI 分区与根分区各自的用途；
  \item 读懂 \texttt{configuration.\allowbreak{}nix} 中常用的选项；
  \item 用 \texttt{nixos-\allowbreak{}rebuild switch} 应用配置，用 \texttt{-{}-rollback} 撤销；
  \item 在配置出错后按确定步骤恢复到可用状态。
\end{itemize}
\end{checkpoint}
